\subsection{Distribution inputs}

This subsection states the distribution estimates used by the rows.
Row S1-0 uses \hyperref[sec:z1-input]{the complete coefficient \textup{(W)}} with
\[
\begin{gathered}
 u=1/500,\quad \Theta(u)=19101/32000,\quad\theta=19101/32000-10^{-4},\\
 B=30,\quad\varepsilon=10^{-4},\quad a=N,\quad K_{\rm scale}=31.
\end{gathered}
\]
Here $0<\sigma\le10^{-10}$ satisfies both restrictions, and
$\tau=\sigma^9$, $Y=N^{\theta\sigma^2}<N^{u/2}$.
We take $C_{\rm box}\ge\lceil2A+32\rceil$ and keep the grid, distinctness masks and smooth-divisor term.
Since $(q,2N)=1$, the shift $a=N$ is admissible.
Then \cite[Theorems~1.1--1.2, Corollary~1.3]{Z1} applies to the complete sum,
and the corollary gives the same Li comparator.
The remaining rows use the following local input.

\begingroup
\renewcommand{\thetheorem}{A}
\begin{theorem}[Local dispersion estimate]\label{thm:distribution-input}
Let $0<\sigma\le10^{-10}$ and let $N$ be large and even.
Put
\[
\begin{gathered}
 \tau=\sigma^9,\quad\delta=30\sigma,\quad\vartheta_{\rm sp}=7/32,
 \quad\eta_L=10^{-4},\\
 \ell_0={53\over412}-{\eta_L\over2}={264897\over2060000},
 \quad c_\sigma={\ell_0\sigma^2\over4},\\
 N/(\log N)^{31}\le X\le N,\quad L=\log X,\quad
 a=N,\quad z_0=X^{1/(\log\log X)^3},\quad u_0=N^{\ell_0\sigma^2}.
\end{gathered}
\]
For an interior Class-II Heath--Brown box $B$, define
\[
\begin{gathered}
 \mathcal S_t(q;a)=1_{t\equiv a\ (q)}-{1_{(t,q)=1}\over\varphi(q)},\quad
 \Delta_B(q;a)=\sum_{\mathbf v\in B}w(\mathbf v)\mathcal S_{\prod v_i}(q;a),\\
 S_B^{\rm PSDT}=\sum_qC_b(q)\Delta_B(q;N).
\end{gathered}
\]
Assume the local conditions below.
Then \cite[Theorem~3.2, (3.13)]{Z1} gives, for fixed $A'>0$,
\begin{equation}\label{eq:distribution-input}
 |S_B^{\rm PSDT}|\ll_{\sigma,31,C_{\rm box},A',\mathcal C_{\rm src}}
 XL^{-A'}+X^{1-c_\sigma}L^C.
\end{equation}
Here $C=C(\sigma,C_{\rm box},A',31,\mathcal C_{\rm src})$ is fixed,
and the constants in $\mathcal C_{\rm src}$ include $C_{\rm ctr},C_{\rm slot}$.
The bound is uniform in boxes, even $N$ and the shift $a=N$.
The constant and threshold may be ineffective.
First $\sigma>0$ is fixed, and then $N\to\infty$.
\end{theorem}
\endgroup

\paragraph{Application checklist.}
Each application uses literal prime sets and the row's own source.
The correspondence with \cite[\S3]{Z1} is
\begin{center}
\begin{tabular}{@{}ll@{}}
\toprule
Condition & Data checked below \\
\midrule
H1 & C1: exact unrestricted coefficient, repeats, $(q,2a)=1$ \\
H2 & C2: labelled factors, last two $u_0$-rough \\
H3 & C2--C3: packing, HB coefficients, majorant, doubled support, LEN \\
H4 & C4: positive shift and its uniform size \\
H5 & C5: selected window \\
H6 & Full-source logarithmic and boundary estimates \\
H7 & Full-family outer coefficient sum \\
\bottomrule
\end{tabular}
\end{center}
\textbf{C1.} For fixed row, prefix cell, AA37 sign and tail vector $b$ we require
\[
\begin{gathered}
 C_b(q)=\sum_{d,\mathbf p}\lambda_d1_{q=d\prod_{j\in J}p_j},\quad |\lambda_d|\le1,\quad p_j\ge u_L\ge u_0,\\
 u_L=L_{\rm rel}^{\sigma^2},\quad q\le X^{0.6},\quad(q,2N)=1.
\end{gathered}
\]
Here $L_{\rm rel}=L_c,L_b,D_c$ for DIRECT, S2, G1.
The labelled primes vary independently over fixed literal sets.
The divisor $d\mid P(u_L;N)$ is squarefree, with the AA37 cap, where
$P(y;N)=\prod_{p<y,\,p\nmid2N}p$.
Selectors are constant on blocks.
The equal-tail, prefix-diagonal, NC, cut and selector-transfer corrections are separated exactly,
and each is estimated separately.

\textbf{C2.} Let $J=J_1\sqcup J_2\sqcup J_3$ be a partition, with $d$ attached to slot 1.
We require, including repeated-prime multiplicities,
\[
 C_b=C_{\rm slot}(\gamma_1*\gamma_2*\gamma_3),\quad
 C_{\rm slot}=\prod_{i=1}^3|J_i|!=O_\sigma(1),\quad |\gamma_i|\le1.
\]
Slots 2 and 3 are $u_0$-rough.
An empty rough slot is $1_{m=1}$.
If $J_1=\varnothing$, put $\gamma_1(m)=\sum_d\lambda_d1_{m=d}$.
The dyadic scales $(Q,R,S)$ belong to slots $(1,3,2)$.
For the test scale $T=T_{\rm II}$ we require
\[
\begin{gathered}
 Q<TX^{-29\sigma},\quad QRS<M_I X^{-89\sigma},\quad
 N_I\ge QX^{\sigma/3},\quad M_I\ge X^{1/2},\\
 Q<TX^{-\sigma},\quad T^2R^2S<X^{1-8\sigma},\quad A_0=X^{1+\sigma/20},\\
 A_0^{\vartheta_{\rm sp}}TQR^2S^5<X^{2-15\sigma},\quad
 (TR/S)^{2\vartheta_{\rm sp}}TQR^2S^5<X^{2-15\sigma}.
\end{gathered}
\]
The labelled-tuple map proves the convolution identity.

\textbf{C3.} The source is either
\[
 1_{\mathbb P}(p)1_{N/(\log N)^{30}<p<(1-\varepsilon_A)N},\quad
 \varepsilon_A=10^{-4},
\]
or the plain product measure
\[
 A_G^{\rm cut}(e)=\sum_{r=3}^4\sum_{\mathbf p\in\mathcal D_r}
 w_r(\mathbf p)1_{e=p_1\cdots p_r}
 1_{\varepsilon_A N<e\le N\{1-(\log N)^{1-C_{\rm box}}\}}.
\]
The ordered $\mathcal D_r$ and bounded $w_r$ are fixed.
Each coordinate is at least $N^{100/1149}$.
All representation multiplicities are retained.
The plain measure is used with its full source predicates.
We apply Heath--Brown's finite identity \cite{ref6} with disjoint boxes of relative width $L^{-C_{\rm box}}$
and bound boundary boxes separately.
The selected $n$ is the eligible variable with largest lower-endpoint exponent.
Ties are broken by variable index, and $m$ is the complement.
For $\zeta\ge z_0$ we require
\[
\begin{gathered}
 \alpha_n\in\left\{\mu(n),1,{\log n\over L}\right\}
 1_{P^-(n)\ge\zeta}1_{I_N}(n),\quad |\alpha_n|\le1,
 \quad|\beta_m|\ll\tau_{23}(m),\\
 I_N\subset(T,2T],\quad N_I\asymp TL^{-C_{\rm box}},\\
 \psi_I(m)=\phi_0((m-m_0)/M_I),\quad M_I\asymp(X/T)L^{-C_{\rm box}},\\
 |m_0|\le C_{\rm ctr}X/T,\qquad 1\le|mn-N|\le X^3.
\end{gathered}
\]
Fixed HB coefficients and tags may be exterior scalars.
The function $\phi_0\in C_c^\infty([-1,1])$ is fixed independently of all parameters,
so its derivative seminorms are fixed and finite.
We require $0\le\phi_0\le1$ and $\phi_0=1$ on $[-1/2,1/2]$.
The complementary interval is majorized by $\psi_I$ on its doubled support.
The shift gap holds throughout the doubled support.
The same integer interval is used, with
\[
 \textbf{LEN:}\qquad N_I\ge(2T)^{1-2^{-126}}.
\]
Tagged comparisons and removed faces have separate estimates.

\textbf{C4--C5.} Every original modulus is coprime to $a=N$, and
\[
 0<a\le X(\log N)^{31}\le A_0,\qquad
 N^{60\sigma}\le T\le N^{1/3+14\sigma},\qquad
 X^{60\sigma}\le T\le X^{1/3+15\sigma}.
\]
For G1, C1--C5 and LEN are checked on the plain $A_G^{\rm cut}$.
This is \textbf{G1-LOCAL}.
Source faces are removed before interior boxes are selected.
The check covers exact HB coefficients, ambient intervals and the fixed majorant.
Lemma~\ref{lem:s2-preparation} gives the direct-source verification,
and Section~4 gives the product-source verification.

\paragraph{Local mechanism.}\label{p2:dispersion-proof}
The argument is \cite[(3.26)--(3.90)]{Z1}.
A first dispersion separates the diagonal and off-diagonal terms.
A second square and Poisson summation give moving-frequency Kloosterman sums.
The spectral estimate and C2 give \eqref{eq:distribution-input}.
The Kloosterman and spectral inputs follow the Deshouillers--Iwaniec
framework and Drappeau's dispersion estimate \cite{ref1,ref3};
the spectral-gap input is Kim--Sarnak \cite{ref7}.
The BFI and Siegel--Walfisz inputs in [Z1] are cited at their sources
\cite{ref2,Siegel1935,Walfisz1936,ref11}.
For S2 the local Mellin data are kept:
\[
\begin{gathered}
 q\asymp Q,\quad r_i\asymp R,\quad s_i\asymp S,\quad n_1,n_2\asymp T,\\
 f=(n_1-n_2)/q\ne0,\quad\eta=\sigma/1000,\quad H_*=40QR^2S^2X^\eta/M_I.
\end{gathered}
\]
On $H_0<|h|\le\min(2H_0,H_*)$, put
$\varsigma_h=\operatorname{sgn}(h)$ and $v=|h|S/(H_0s_1)$.
Let $w$ be smooth, supported in $(1/4,4)$ and equal to 1 on $[1/2,2]$.
Then
\begin{equation}\label{p2:eq:dispersion-mellin}
\begin{gathered}
 W_Y(v)=w(v)\widehat\phi_0(Y_2v)e(Y_1v),\quad
 Y_1={\varsigma_hH_0(a/n_1-m_0)\over Sqr_1r_2s_2},\quad
 Y_2={\varsigma_hH_0M_I\over Sqr_1r_2s_2},\\
 |Y_1|+|Y_2|\ll L^{C_{\rm box}+31}X^\eta=:L_{\rm osc},\qquad
 \int_{\mathbb R}\sup_Y|\widetilde W_Y(it)|\,dt\ll L_{\rm osc}.
\end{gathered}
\end{equation}
Both nonmodular phases remain in $W_Y$.
Put $m(t)=\sup_Y|\widetilde W_Y(it)|$.
We normalize by $m(t)>0$ and restore it before integration;
for $m(t)=0$ the integrand vanishes.
The details are in \cite[(3.56)--(3.65)]{Z1} and \S\ref{sec:s2-3}.

Type I uses \cite[Lemmas~3.3--3.5, \S3.11]{Z1} only after preparation.
Lemma~\ref{lem:s2-preparation} verifies its source and scale hypotheses.
The product source has only Class-II interior boxes (Section~4).
The respective choices are
\[
 D_f=X^{1/20},\quad V\ge qX^{1/10};\qquad D_f=X^{7.5\sigma/2}.
\]
At these prepared scales \cite[(3.102)--(3.115)]{Z1} applies.

\paragraph{Global conditions.}
The two global conditions bound the outer coefficients and the logarithmic and boundary errors.
For H7 we sum every outer coefficient:
{\scriptsize
\[
 \begin{gathered}
 \sum_j|c_j|\ll_{\sigma,C_{\rm box},\mathcal C_{\rm src}}
 (\log N)^{K_0},\\
 K_0\le24C_{\rm box}+26,\qquad A'=A+24C_{\rm box}+31.
 \end{gathered}
\]
}
Then \cite[(3.92)]{Z1} gives
\[
 E_{\rm II}\ll NL^{K_0-A'}+N^{1-c_\sigma}L^{K_0+C}.
\]
Since $K_0-A'\le-A-5$, the first term is $\ll NL^{-A-5}$.
For fixed $\sigma>0$, the second is also $O(NL^{-A-5})$.
The numerical application uses the per-box precision in (1.28).
For H6 we sum all labels, shells, boxes and representation multiplicities.
We require fixed $c_{\log},c_R$ with
\[
\begin{gathered}
 E_{\rm log}\ll N(\log N)^{c_{\log}-C_{\rm box}},\quad
 E_{\rm bdry}\ll N(\log N)^{c_R-C_{\rm box}/2},\\
 C_{\rm box}\ge\max\{A+c_{\log},2(A+c_R)\}.
\end{gathered}
\]
For the prime source, $(c_{\log},c_R)=(18,16)$
\cite[\S3 and Appendix~B]{Z1}.
Section~4 verifies the product-source values separately.
Together these give the aggregate $O(N(\log N)^{-A})$ error.

The next lemma checks packing on these rows.

\begin{lemma}[Packing on the actual row domains]\label{thm:packing-input}
Consider a row in the table below, including the three empty prefixes.
Let $c$ be its retained cell, at $\Theta_c=\inf_c\widetilde H$,
with its literal source domain and block selector.
For a nonempty prefix, let $t_1\ge\cdots\ge t_r$, $1\le r\le3$.
For either AA37 sign $\nu$, let
$\mathbf D\in\mathcal D^\nu(L_c)$ be any finite tail of length $k\ge0$.
Put
\[
\begin{gathered}
 \ell_c=\Theta_c-\eta_L-\sum_{i\le r}t_i\ge\ell_0,\quad L_c=N^{\ell_c},
 \quad z\le L_c\ (\text{upper}),\quad z^2\le L_c\ (\text{lower}),\\
 \rho=1+\sigma^9,\quad \delta=30\sigma,\quad\vartheta_{\rm sp}=7/32,
 \quad n\in[60\sigma,1/3+14\sigma],\\
 T=N_{\rm II}=N^n,\quad h=\rho(\Theta_c-\eta_L)+2\delta,
 \quad g=n-\delta,\quad A_{\max}=X^{1+\sigma/20}.
\end{gathered}
\]
Sort the endpoint exponents as $e_1\ge\cdots\ge e_m$, $m=r+k$.
Repeats are kept as distinct labelled atoms, and the list is padded with zeros.
Assume the DIRECT or plain G1 intervals satisfy
\[
 N_I\gg T(\log X)^{-C_{\rm box}},\qquad
 \log_NM_I\ge2/3-14.02\sigma,\qquad M_I\ge X^{1/2}.
\tag{P6}
\]
These bounds are uniform in row labels and boxes.
Then a deterministic labelled partition satisfies C2 at the same $\Theta_c$.
It attaches $d_{AA}$ to slot 1 and allows unit slots.
After dyadic and base transfer, the four spectral surpluses are at least
\[
 \left({7271\over256},\ 52-{1\over100},\
 {95221\over1280}-{1\over50},\ {22407\over256}-{1\over50}\right)\sigma
 >(28,51,74,87)\sigma.
\]
All four lower-scale conditions hold, with $\Delta D_{\rm PACK}=0$.
For G1-3 the cutoff is $z=N^{t_3}$ after the PR peel.
\end{lemma}

\begin{proof}
Put
\[
\begin{gathered}
 u=1/500,\quad\alpha=100/1149,\quad\beta=50/309,\quad\mu=25/128,\quad\lambda=93/400,\\
 t_*=7109/28000,\quad t_{\rm dd}=25/[32(1+3\rho)],\quad\overline H=153/256,\quad T_0=509/1000-\eta_L.
\end{gathered}
\]
Here $v,w,H,\widetilde H$ are as in \eqref{p2:eq:row-level-function}.
Each interval $[N^{t_i},N^{\rho t_i})$ is clipped to its row domain and keeps its label.
Missing coordinates lie in $[0,w_z]$, where $z=N^{w_z}$.
For $f_1\ge\cdots\ge f_m\ge0$, write
\[
 \mathrm{WF}(y):\quad\sum_{i<j}f_i+2f_j\le y\quad(1\le j\le m).
\tag{P1}
\]
The proof follows \cite[Theorem~4.1, (4.1)--(4.5)]{Z1},
with auxiliary identities \cite[Lemmas~4.2--4.4]{Z1}
and two-bin step \cite[Lemma~4.5]{Z1}.
The two-bin rule retains labels, repeats and zero padding.
Put $d_*=h-2\delta$ and $\Sigma_s=\sum_{i\in J_s}e_i$.
We verify the hypotheses of that theorem on these actual rows.

By \cite[Lemma~4.2]{Z1}, $v$ is concave and $\widetilde H\le\overline H$.
For $x_1\in[1/10,\rho t_*]$,
\[
 v(x_1)\ge\min\{v(1/10),v(127/500)\}>509/1000.
\]
For $x_1<1/10$, the active $w$ branch gives
\[
 \min\{(132-68x_3)/214,(153-96x_3)/256,1-2x_2\}
 \ge717/1280\ge11/20\ge(1+x_1)/2.
\]
Hence
\[
 509/1000\le\Theta_c\le\overline H,\qquad
 \Theta_c\ge m(t_1):=\min\{(1+t_1)/2,57/64-3\rho t_1/2\}.
\tag{P5}
\]
The first branch is smaller for $t_1\le t_{\rm dd}$.
At rank 1, $2t_1<2t_*=7109/14000<T_0$.
The remaining prefix checks are
\[
 \sum_{i<j}t_i+2t_j<\Theta_c-\eta_L\quad(1\le j\le r).
\tag{P4}
\]
For S1, S2, G1 the cutoff is $z=N^u$ at $r\le2$.
At $r=3$ it is $z=N^{t_3}$, after PR for G1.
The other rows use $z=N^\alpha$.

\begin{longtable}{@{}p{0.29\textwidth}p{0.66\textwidth}@{}}
\toprule
Rows and literal domains & Prefix checks and selectors \\
\midrule\endhead
S1-1, S1-2, S1-3; $[u,\alpha)$
& $(j+1)t_1\le4\alpha=400/1149<T_0$. \\[3pt]
S2-1, S2-2, S2-3; $[u,\beta)$
& Rank 2: $t_1+2t_2\le3\beta=150/309<T_0$.
Rank 3: $t_1<\beta<t_{\rm dd}$ and the whole-$K$ selector
$t_1+2t_2+(4+2\tau)t_3\le249/250+2\tau\beta$.
Thus $t_1+2t_2+4t_3<1-2\eta_L$;
$1-2\eta_L-249/250=19/5000>2\tau\beta$. \\[3pt]
G1-1, G1-2, G1-3; $[u,\mu)$
& Rank 2: $t_1+2t_2\le3t_1$,
$\mu<(1-2\eta_L)/5$, and
$3/256-\eta_L-(3/2)\tau\mu>0$.
Rank 3 requires $\Theta_c-\eta_L-t_1-t_2-(3+2\tau)t_3\ge0$.
Then $t_1+t_2+2t_3\le\Theta_c-\eta_L-(1+2\tau)t_3$,
$t_3\ge u/\rho>0$. \\[3pt]
S3-A: $[\alpha,\mu)$;
S4: $[\mu,1/4)$;
S5-low: $[1/4,t_*)$;
S7: $[\mu,\lambda)$
& One prefix each; $2t_1<2t_*<T_0$. \\[3pt]
S8-A: $\alpha\le t_2<t_1<\beta$
& $t_1+2t_2\le3\beta=150/309<T_0$. \\[3pt]
S9-A: $\beta\le t_1<\mu$, $\alpha\le t_2<\beta$
& $t_1<1-4\beta-2\eta_L$ for the first branch;
$3886931/49440000-(3/2)\tau\mu>0$ for the second. \\[3pt]
S10-Ret (the retained S10 row): $\mu\le t_1<\lambda$, $\alpha\le t_2<\beta$
& $t_2\le\operatorname{sel}(t_1)$,
$C_0=57/64-u-(11/5)\alpha$,
$\operatorname{sel}(t)\le B_\tau(t)\le\max\{\alpha,C_0-(5/2+3\tau/2)t\}$.
For $\Theta_c-u-t_1-t_2\ge(11/5)\alpha$,
the second branch in (P5) is active, since $t_1\ge\mu\ge t_{\rm dd}$.
Its affine alternative gives equality;
the $\alpha$ alternative gives $265391/9192000-(3/2)\tau\lambda>0$.
Thus $t_1+2t_2\le\Theta_c-u-(11/5)\alpha+t_2<\Theta_c-\eta_L$,
with binding prefix-mask gap
$u+(11/5)\alpha-\eta_L-\beta=12449531/394490000>0.03155$. \\[3pt]
S1-0, S2-0, G1-0; empty prefix
& $\Theta=H(u,u,u)=19101/32000$; tail endpoints $\le u$.
The active $w$ branch gives
$\widetilde H(\bar e_1,\bar e_2,\bar e_3)\ge
\min\{(132-68u)/214,(153-96u)/256,1-2u\}=19101/32000$. \\
\bottomrule
\end{longtable}

A prefix endpoint lies in $[t_i,\rho t_i]$.
A tail endpoint satisfies $\epsilon'_j\le\min\{\rho s_j,w_z\}$,
from $[N^{s_j},\min\{N^{\rho s_j},z\})$.
Every prefix prime is at least $z$.
The AA37 masks give
\[
 \sum_{j\le2l}s_j+3s_{2l+1}<\ell_c\ (\nu=+),\quad
 \sum_{j\le2l-1}s_j+3s_{2l}<\ell_c\ (\nu=-),\quad 2s_1<\ell_c.
\]
We drop one copy at cubic ranks and use $s_{j-1}\ge s_j$ between them.
Then $\sum_{i<j}s_i+2s_j<\ell_c$ at every tail rank.
The prefix sum is then added at these ranks.
If $s_1,\ldots,s_q$ precede $t_i$, then $t_i\le s_q$ and
\[
 \sum_{a<i}t_a+\sum_{b\le q}s_b+2t_i
 \le\sum_{a\le i}t_a+\left(\sum_{b<q}s_b+2s_q\right)
 <\Theta_c-\eta_L.
\]
Endpoint inflation by $\rho$ gives $\mathrm{WF}(d_*)$.
The first $r$ sorted endpoints lie in $[t_j,\rho t_j]$.
The rest are tails or zero pads in $[0,w_z]$.
Their first three entries therefore lie in the closed parent box.
Hence
\[
\begin{gathered}
 \Theta_c\le\widetilde H(\bar e_1,\bar e_2,\bar e_3),\quad
 \Theta_c-h=\eta_L-\tau(\Theta_c-\eta_L)-60\sigma>0.97\eta_L,\\
 1/2<h\le\overline H,\quad 0<\delta\le3\cdot10^{-9},\quad
 g=n-\delta\in[\delta,1/3-\delta/2].
\end{gathered}
\]
These verify \cite[Theorem~4.1, (4.2)]{Z1} for every finite tail.
The mixed-prefix checks above verify \cite[Theorem~4.1]{Z1} directly.
The pure-tail result \cite[Lemma~4.7]{Z1} is not used.
For the empty prefixes, the table gives the same hypotheses.

We attach $d_{AA}<L_c^\sigma=N^{\sigma\ell_c}$ to slot 1.
Then
\[
\begin{gathered}
 \log_NQ\le\Sigma_1+\sigma\overline H,\quad\log_NS\le\Sigma_2,\quad\log_NR\le\Sigma_3,\\
 n=g+\delta,\quad\log_NA_{\max}\le1+\sigma/20.
\end{gathered}
\]
By \cite[(4.3)]{Z1}, the four excess exponents are
\[
\begin{aligned}
 -(30-\overline H)\sigma&=-(1+7271/256)\sigma,\qquad -2\delta=-60\sigma,\\
 -3\delta+(\vartheta_{\rm sp}/20+\overline H)\sigma
 &=-(15+95221/1280)\sigma,\\
 -(3+2\vartheta_{\rm sp})\delta+\overline H\sigma
 &=-(15+22407/256)\sigma.
\end{aligned}
\]
Dyadic representatives within a factor two cost $24\log2/\log N<\sigma/1000$.
For $k_X=\log_XN=1/(1-\omega)$, $\omega=\log_N(N/X)$,
$0\le k_X-1<0.01\sigma$.
The first exponent improves.
The second costs at most $0.01\sigma$, and the last two cost at most $0.02\sigma$.
These give the stated spectral margins.

The factorization follows \cite[Lemma~4.8, (4.12)--(4.13)]{Z1}.
The small-prime divisor is unique in slot 1.
Each labelled slot has at most $|J_s|!$ multiset orderings.
Thus its finite-sum identity holds for these literal sets, including repeats.
Normalization gives $|\gamma_s|\le1$; C2 covers empty slots.
Lemma~4.8 bounds tail-collision support only; source corrections remain separate.
Also,
\[
 L_c^{\sigma^2}\ge u_0,\quad p_{\rm prefix}\ge N^u,\quad
 s_j\ge\sigma^2\ell_c,\quad\sum_js_j<\ell_c,\quad
 k<\sigma^{-2},\quad C_{\rm slot}=O_\sigma(1).
\]
Hence slots 2 and 3 are $u_0$-rough.

Condition (P6) is verified separately for DIRECT and plain G1.
The first margin gives $Q<TX^{-29\sigma}$, so
\[
 {N_I\over QX^{\sigma/3}}\gg
 (\log X)^{-C_{\rm box}}X^{(29-1/3)\sigma}\longrightarrow\infty.
\]
The last WF mask gives $\sum_s\Sigma_s\le d_*<\overline H$.
Thus
\[
\begin{gathered}
 \log_N(QRS)\le153/256+\sigma,\\
 \log_NM_I-\log_N(QRS)\ge
 2/3-153/256-15.02\sigma=53/768-15.02\sigma>89\sigma.
\end{gathered}
\]
Since $\sigma\le10^{-10}$, $QRS<M_I X^{-89\sigma}$.
Condition (P6) also gives $M_I\ge X^{1/2}$.
The same rule covers zero pads and unit slots.
For G1-3, the comparison $N^{\rho t_3}\to N^{t_3}$ is proved separately.
Packing uses the latter cutoff at the original $\Theta_c$,
so its level loss is zero.
\end{proof}

Sections~2--3 verify direct-row source conditions; Section~4 verifies plain G1.
Each source comparison and removed boundary term has its separate estimate.
